\documentclass[11pt,titlepage]{article}


\usepackage{packages}


\title{The 2D Schwinger Model\\On Lattice}
\author{Jenish Adhikari}
\date{\today}

\begin{document}

\maketitle
\tableofcontents
\newpage
\section{Introduction}
\label{sec:intro}
Lattice Quantum Chromodynamics (LQCD) provides a non-perturbative framework for studying strongly interacting gauge 
theories from first principles. By discretizing space-time onto a finite lattice, quantum field theories can be formulated 
in a way that makes them amenable to numerical Monte Carlo simulations. This approach has become one of the most powerful tools 
for investigating phenomena such as confinement, mass generation, and chiral symmetry breaking.

As an introduction to lattice gauge theory methods, this project studies the two-dimensional Schwinger model, 
which is Quantum Electrodynamics (QED) in 1+1 dimensions with gauge group $U(1)U(1)$. Although significantly simpler than 
4-D QCD, the Schwinger model shares many important qualitative features with non-abelian gauge theories.
In particular, it exhibits confinement, topological sectors, etc. One major advantage of the Schwinger model is its reduced dimensionality. 
In 1+1 dimensions, gauge degrees of freedom are highly constrained by Gauss’ law, simplifying both the analytical structure of the theory and 
its numerical implementation. Despite this simplicity, the model retains essential conceptual ingredients of LQCD, including gauge invariance, 
compact link variables, fermion discretization issues, and the need for importance sampling algorithms such as Hybrid/Hamiltion Monte Carlo (HMC).

In this project, the continuum Schwinger model is formulated on a Euclidean space-time lattice using compact $U(1)$ gauge links. 
The Wilson gauge action is employed, and the path integral is evaluated numerically using Monte Carlo techniques. Key observables such as 
the plaquette expectation value and topological charge are measured to study equilibration, confinement properties, etc.

The methodology developed here closely parallels that used in full LQCD simulations. Although the gauge group is abelian and the dimensionality 
is reduced, the essential computational framework, gauge action construction, fermion integration, and Markov-chain 
sampling can be directly generalized to $SU(N)$ gauge theories, including realistic QCD. Thus, this project serves as a conceptual introduction 
to lattice field theory. This report begins with a brief overview of the theoretical background of the field theory, followed by a description of 
its lattice formulation. The simulation methodology, including the implementation of the Hybrid Monte Carlo algorithm, is then presented. 
Finally, the procedures for data analysis and the results obtained from the numerical simulations are discussed.



\section{Theory}
\subsection{QFT}
Let's consider the quantum field theory of QED, the langragian density (later simply langragian) is givien as \cref{eqn:qed_full}.

\begin{equation}\label{eqn:qed_full}
    \mathcal{L}\big(\bar{\psi},\psi,A_\mu\big)
    =
    \sum_{f=1}^{N_f}
    \bar{\psi}^{(f)}(x)\Big[\gamma^\mu i \big(\partial_\mu + i g A_\mu(x)\big) - m\Big]\psi^{(f)}(x)
    -\frac{1}{4} F_{\mu\nu}(x)F^{\mu\nu}(x)
\end{equation}
In Path Integral formalisim an expectation value is given by
\begin{equation}
    \langle O \rangle = \frac{1}{Z} \int \mathcal{D}A_\mu\mathcal{D}\psi\mathcal{D}\bar{\psi} O[A_\mu,\psi,\bar{\psi}] e^{iS}
\end{equation}
\begin{equation}
    Z = \int \mathcal{D}A_\mu\mathcal{D}\psi\mathcal{D}\bar{\psi} e^{iS}
\end{equation}
After permofing a Wick's rotation \cref{eqn:wick_rot} the action becomes \cref{eqn:action_E}\footnote{Note the gamma matrices changes after rotation.}.

\begin{align}
    S &= \int d^4x \mathcal {L} \\
    t &\to -i\tau \label{eqn:wick_rot}\\
    Z &= \int \mathcal{D}A_\mu\mathcal{D}\psi\mathcal{D}\bar{\psi} e^{-S_E} \\
    S_E &= \int d^4x  
    \bar{\psi}(x)\Big[\gamma^\mu\big(\partial_\mu + i g A_\mu(x)\big) + m\Big]\psi + \frac{1}{4} F_{\mu\nu}(x)F^{\mu\nu}(x) \label{eqn:action_E}\\
    \gamma_E^i &= i \gamma_M^i
\end{align}
\subsection{Lattice discretization}
Space-time is discretized into a two-dimensional lattice with spacing $a$. 
Gauge fields are represented by compact link variables

\begin{equation}
U_\mu(n) = e^{iA\mu(n)} \in U(1)
\end{equation}
defined on links between lattice sites $n$ and $n + \mu$.
For 2D-Schwinger model the fermion matrix read as,

\begin{equation}
     M[n',n]^{\alpha\beta} = (m_0+2)\delta^{\alpha\beta}\delta_{n',n} -\frac{1}{2}\sum\limits_{\mu=0}^{1}\{(1-\sigma_\mu)^{\alpha\beta}U_{\mu}(n')\delta_{n'+\mu,n} +
        (1+\sigma_\mu)^{\alpha\beta}U^{\dagger}_{\mu}(n'-\mu)\delta_{n'-\mu,n}  \},
\end{equation}
where $\sigma$ are Pauli sigma matrix. Now the action can be sperated in purge gauge part and fermion part as \cref{eqn:disc_action}. 
In limit $a \to 0$ the continuum field langragian is obtained~\cite{gupta1998introduction}. 
\begin{align}\label{eqn:disc_action}
    S_G &= \beta\sum_{n\in \Lambda}\sum_{\mu<\nu} \mathcal{R}Tr(1-P_{\mu\nu}(n))\\
    S_F &= \bar{\psi} M \psi
\end{align}

Consider the path integral, using properties of Grassmann number \cref{eq:grassmann_rules} the fermion field can be integrated out to \cref{eqn:only_gaugePart}. 
In numerical methods we want to sample from the distribution, computing the determinant scales complexity with  $\mathcal{O}(n!)$ and space with
$\mathcal{O}(n^2)$, where $n$ is the volume of the lattice. We can implement $M^\dagger M \chi$ as a vector and solve linear equation to get inverse
as $(M^\dagger M)^{-1} \eta$ with $ \eta = M^\dagger \chi$, where $\chi \sim \mathcal{N}(0,2)$
\begin{align}
    \langle \mathcal{O} \rangle &= \frac{1}{Z} \int \mathcal{D}A_\mu\mathcal{D}\psi\mathcal{D}\bar{\psi} \mathcal{O}[A_\mu,\psi,\bar{\psi}] e^{-S_G - S_F}\\
     &= \frac{1}{Z} \int \mathcal{D}A_\mu \mathcal{O}[A_\mu] \det(M^\dagger M) e^{-S_G} \label{eqn:only_gaugePart}\\
     &= \frac{1}{Z} \int \mathcal{D}A_\mu\mathcal{D}\eta\dagger \mathcal{D} \eta \mathcal{O}[A_\mu]  e^{- \left(S_G + \eta\dagger (M^\dagger M)^{-1} \eta \right)} \label{eqn:inverse_f_M}
\end{align}

\section{Simulation}
We want to sample from the distribution $e^{-S_E}$. A local Metropolis update is extremely expensive 
due to dynamical fermion matrix and suffers from larger autocorrelation. To generate gauge configurations of desired distributed, 
we use the Hybrid Monte Carlo algorithm and implement global Metropolis update.

\subsection{Hamiltonian Formulation}
Once we add a fictitious conjugate momenta $P_\mu(x)$ for each gauge link, $S_E \to H$. Now, we have a Gaussian
integral, which can be calculated analytically. Thus the addition of momentum doesn't affect the expectation value, due to same contirbution 
in numerator and denominator.
\begin{equation}
H = \sum_{x,\mu} \frac{1}{2} P_\mu(x)^2 
+ S_{\text{eff}}[U].
\end{equation}

The partition function becomes
\begin{equation}
Z = \int \mathcal{D}U\, \mathcal{D}P \;
e^{-H[U,P]}.
\end{equation}

\subsection{Molecular Dynamics}

To propose a new gauge link, the fields are evolved in fictitious simulation time $\tau$
according to Hamilton's equation of motion:
\begin{align}
\frac{dU_\mu(x)}{d\tau} &= i P_\mu(x) U_\mu(x), \\
\frac{dP_\mu(x)}{d\tau} &= 
- \frac{\partial S_{\text{eff}}}{\partial \theta_\mu(x)},
\end{align}

The total force can be splited into
\begin{equation}
- \frac{\partial S_{\text{eff}}}{\partial \theta_\mu(x)}= F = F_g + F_f,
\end{equation}
where $F_G$ is the gauge force and $F_F$ is the fermionic force.

\begin{align}
    F_G = -\beta \Im[U_\mu(n)K_\mu(n)] &   \\
\textrm{with } K_\mu(n) = U_\nu(n+\mu)U_\mu^\dagger(n+\nu)U_\nu^\dagger(n) &+ U_\nu^\dagger(n+\mu-\nu) U_\mu^\dagger(n-\nu)U_\nu(n-\nu)|_{\mu \neq \nu } \\
F_F = \Im[\chi^\dagger(n)(1-\sigma_\mu)U_\mu(n)\xi(n+\mu) &- \chi^\dagger(n+\mu)(1+ \sigma_\mu)U_\mu^\dagger(n)\xi(n)]\\
\textrm{with } \xi = (M^\dagger M)^{-1}\Phi \textrm{ and } \chi = M \xi
\end{align}

The molecular dynamics equations are solved numerically
using the leapfrog integrator, which is symplectic,
time-reversible, and area-preserving. The integrator introduces $\mathcal{O}(\epsilon^2)$ errors,
which are corrected by the Metropolis step.

\subsection{Hamiltion Monte Carlo}

After a trajectory, compute the new Hamiltonian $H_{\text{new}}$.
The configuration is accepted with probability
\begin{equation}
P_{\text{acc}} =
\min\left(1, e^{-(H_{\text{new}} - H_{\text{old}})} \right).
\end{equation}

\section{Analysis}
\begin{figure}[H]
    \centering
    \includegraphics[width = \textwidth]{runnig_mean_4x4.png}
    \caption{Running mean of Plaquette for 4x4-Lattice.}
    \label{fig:run_mean_4_L}
\end{figure}

The equilibration of the simulated system is a crucial first step in numerical lattice simulations. 
\Cref{fig:run_mean_4_L} shows the running mean of the plaquette.
The plaquette $P_{\mu\nu}(n)$ measures the product of link variables around an elementary square on the lattice:
\begin{equation}
    P_{\mu\nu}(n) = \text{Re}[U_\mu(n)U_\nu(n+\mu)U_\mu^\dagger(n+\nu)U_\nu^\dagger(n)]
\end{equation}

The running mean demonstrates that the system rapidly converges to its equilibrium value within the first few hundred trajectories. 
The initial transient behavior, where the plaquette value drifts from its starting configuration, indicates that the system is thermalizing. 
Once the running mean stabilizes, the system has reached thermal equilibrium and subsequent measurements represent independent samples from the desired distribution. 
Similar observation can be seen for different lattice size \cref{fig:run_mean_5_L} and \cref{fig:run_mean_6_L}

\subsection{Topological Charge Evolution}
The integer-valued topological charge functional for 2D- Schwinger model is given by~\cite{elser1997topological},
\begin{equation}
Q\;=\; \frac{1}{2\pi}\int d^{2}x\;F_{01}(x)\,.
\end{equation}
One can measure the discreet topological charge with
\begin{equation} \label{eqn:discreet_Q}
Q=\frac{1}{2\pi}\sum_{n \in \Lambda}\arg P_{\mu\nu}(n)
\end{equation}
\begin{figure}[H]
    \centering
    \includegraphics[width=\textwidth]{Q_4.png}
    \caption{Topological charge for 4x4-Lattice}
    \label{fig:Q_4_L}
\end{figure}

\Cref{fig:Q_4_L}, \cref{fig:Q_5_L} and \cref{fig:Q_6_L} displays the distribution of the topological charge from the HMC simulation. 
One can observe that, the topological charge fluctuates between different integer values, 
demonstrating that the HMC algorithm successfully tunnels between distinct topological sectors.

This behavior is highly desirable and contrasts sharply with some local update algorithms that become 
frozen or have very low tunnelling probability in a single topological sector~\cite{elser1997topological}. 
The ability to explore multiple topological sectors ensures that:
\begin{itemize}
    \item The simulation samples from the complete configuration space
    \item Physical observables are not biased by topological sector restrictions
\end{itemize}

\subsection{Autocorrelation Analysis}

\begin{figure}
    \hspace*{-4cm}
    \includegraphics[width=1.6\textwidth]{autocorr_5_Lattice.png}
    \caption{Autocorrelation of topological charge (left) and Wilson loop W(R,T) for R = 1, T = 1 in 5x5-Lattice (right)}
\end{figure}

One of the primary advantages of the Hybrid Monte Carlo algorithm over local update methods is the significantly reduced autocorrelation time. 
\Cref{fig:Q_4_L} shows normalized autocorrelation function $\Gamma$ for topological charge  on left  and Wilson loop $W(R,T)$  on right. Both shows 
significant low integrated autocorrelation time. For topological charge $\tau_{\text{int}} \sim 1.28$ and for Wilson loop $W(1,1)$ $\tau_{\text{int}} \sim 0.5$ 
for 5x5-Lattice. Both the topological charge and the Wilson loop demonstrate
the similar autocorrelation and integrated autocorrelation time for different lattice size.

\subsection{Confinement Evidence}

\begin{figure}
    \centering
    \includegraphics[width=\textwidth]{pot_V_T_1.png}
    \caption{Potential plotted as -ln(W(R,T)) with T = 1. The potential growes linearly showing confinement.}
    \label{fig:pot_V_T_1}
\end{figure}

A defining characteristic of QCD and 2D-Schwinger model is the phenomenon of confinement. 
The confinement in Schwinger model is due to Gauss' law in 1-Spatial dimension $\partial_x E \sim q$.
Unlike free theories where the potential between charges grows logarithmically (or vanishes), confining theories exhibit a linear potential at large distances. 
In \cref{fig:pot_V_T_1} a linear fit to \cref{eq:V_R} demonstrates this confinement property by plotting the effective potential,

\begin{equation}
    \label{eq:V_R}
    V(R) = -\ln W(R,T) / T
\end{equation}

where $W(R,T)$ is the Wilson loop operator for a rectangular closed loop with spatial extent $R$ and temporal extent $T$~\cite{lecturenote}. 
The linear growth of the potential with separation distance $R$ shows confinement predicted by the theory.

This linear confinement behavior emerges naturally from the HMC simulations without any external constraint and is a consequence of the gauge-invariant 
dynamics of the Schwinger model. The successful observation of confinement demonstrates that:
\begin{itemize}
    \item The lattice formulation correctly preserves the essential physics
    \item The HMC algorithm efficiently samples the relevant gauge configurations
    \item The numerical implementation accurately reproduces theoretical predictions
\end{itemize}

The confinement phenomenon observed in these simulations validates the approach and confirms that the 2D Schwinger model, 
studied via lattice Monte Carlo methods, provides valuable insights for understanding gauge theories.

\newpage

\section{Conclusions}

The numerical simulations of the 2D Schwinger model using the Hybrid Monte Carlo algorithm have successfully demonstrated several key physical and computational results:

\begin{enumerate}
    \item \textbf{Rapid Equilibration:} The system reaches thermal equilibrium within a few hundred HMC trajectories, as evidenced by the stabilization of the plaquette running mean. This rapid convergence validates the choice of HMC for exploring the configuration space efficiently.
    
    \item \textbf{Topological Sector Exploration:} The algorithm successfully tunnels between different topological sectors, with the topological charge fluctuating between integer values throughout the simulation. This demonstrates that HMC provides unbiased sampling across the entire configuration manifold, avoiding the freezing that plagues local update algorithms.
    
    \item \textbf{Reduced Autocorrelation:} Both the topological charge and Wilson loop observables exhibit short autocorrelation times, indicating statistically independent measurements. This efficiency gain over local methods translates to faster convergence of statistical estimates and reduced computational cost for achieving target precision.
    
    \item \textbf{Confinement Behavior:} The extracted interquark potential demonstrates clear linear growth with distance, confirming the confinement property expected from Schwinger model. 
\end{enumerate}


Despite the successful demonstration of HMC's capabilities, the simulations also reveal important limitations that become increasingly significant as the lattice volume increases.
The leapfrog integrator, introduces systematic errors that scale as $\mathcal{O}(\epsilon^2)$ per step. 
For a fixed trajectory length $\tau$, the leapfrog (second-order) integrator
induces a Hamiltonian violation whose typical size scales as
\begin{equation}
\Delta H_{\mathrm{rms}} \sim \sqrt{V}\,\epsilon^{2},
\qquad
\epsilon = \frac{\tau}{N_{\text{step}}}.
\end{equation}
Consequently, increasing the lattice volume $V$ leads to larger integration
errors and a reduced Metropolis acceptance probability.
To maintain an acceptance rate above $50\%$, the step size $\epsilon$ must be
decreased, which in practice requires increasing the number of leapfrog steps
$N_{\text{step}}$ at fixed $\tau$.
This significantly increases the computational cost, since each additional
molecular-dynamics step requires an evaluation of the  force.







\newpage
\printbibliography

\begin{equation}
\int d\psi = 0,
\qquad
\int d\bar{\psi} = 0,
\qquad
\int d\psi \, \psi = i,
\qquad
\int d\bar{\psi} \, \bar{\psi} = i,
\end{equation}

\begin{equation}
\int d\bar{\psi} \, d\psi = 0,
\int d\bar{\psi} \, d\psi \; \psi = 0,
\int d\bar{\psi} \, d\psi \; \bar{\psi} = 0,
\int d\bar{\psi} \, d\psi \; \bar{\psi} \psi = 1.\label{eq:grassmann_rules}
\end{equation}

\begin{figure}
    \centering
   \begin{subfigure}{0.45\textwidth}
        \includegraphics[width=\textwidth]{runnig_mean_5x5.png}
        \caption{Running mean of Plaquette for 5x5-Lattice.}
        \label{fig:run_mean_5_L}
   \end{subfigure}
   \begin{subfigure}{0.45\textwidth}
        \includegraphics[width=\textwidth]{runnig_mean_6x6.png}
        \caption{Running mean of Plaquette for 6x6-Lattice.}
        \label{fig:run_mean_6_L}
   \end{subfigure}
   
\end{figure}

\begin{figure}
    \centering
    \begin{subfigure}{0.45\textwidth}
        \includegraphics[width=\textwidth]{Q_5.png}
        \caption{Topological charge for 5x5-Lattice}
        \label{fig:Q_5_L}
    \end{subfigure}
     \begin{subfigure}{0.45\textwidth}
        \includegraphics[width=\textwidth]{Q_6.png}
        \caption{Topological charge for 6x6-Lattice}
        \label{fig:Q_6_L}
    \end{subfigure}


\end{figure}
\end{document}